\section{From the line metaphor to a research object}
\label{sec:tube}
Along a trajectory, the PRM produces the line
$\hat{\bm r}(\tau)=(\hat r(h_1),\ldots,\hat r(h_T))$. A tube should contain
plausible target rewards, not merely decorate that line with arbitrary error bars.
Define lower and upper envelopes $L(h),U(h)$, a domain $\mathcal H_B$ indexed by
an optimization budget $B$, and a set of functions
\[
 \U_B=\{r:\mathcal H_B\to[0,1]:L(h)\le r(h)\le U(h),\quad
 r\text{ satisfies specified coupling constraints}\}.
\]
The coupling says which combinations of step errors are simultaneously plausible.
Independent intervals form a box; shared parameters form a correlated tube.
This distinction is the main geometric research question.

\begin{figure}[ht]
\centering
\begin{tikzpicture}[x=0.92cm,y=4cm]
\draw[-{Stealth}] (0,0) -- (7.3,0) node[right]{step};
\draw[-{Stealth}] (0,0) -- (0,1.03) node[above]{reward};
\fill[blue!12] (1,.88)--(2,.94)--(3,.84)--(4,.95)--(5,.98)--(6,.92)--(7,.87)
 --(7,.35)--(6,.18)--(5,.28)--(4,.62)--(3,.56)--(2,.73)--(1,.72)--cycle;
\draw[blue,dashed] (1,.88)--(2,.94)--(3,.84)--(4,.95)--(5,.98)--(6,.92)--(7,.87);
\draw[blue,dashed] (1,.72)--(2,.73)--(3,.56)--(4,.62)--(5,.28)--(6,.18)--(7,.35);
\draw[blue,thick] (1,.80)--(2,.84)--(3,.70)--(4,.79)--(5,.87)--(6,.78)--(7,.73);
\draw[orange!85!black,thick] (1,.79)--(2,.81)--(3,.65)--(4,.74)--(5,.43)--(6,.30)--(7,.52);
\node[blue,anchor=west] at (7,.86) {$U$};
\node[blue,anchor=west] at (7,.72) {$\hat r$};
\node[blue,anchor=west] at (7,.34) {$L$};
\node[orange!85!black,anchor=west] at (7,.51) {$r^*$};
\foreach \x in {1,...,7}{\draw (\x,0.01)--(\x,-.01) node[below]{\x};}
\end{tikzpicture}
\caption{Illustrative tube with heterogeneous width. All values are schematic,
not measured results. Drawing a target line inside the band does not establish
statistical coverage or describe which correlated functions are admissible.}
\end{figure}

\subsection{The design specification}
Before implementing any method, specify:
\begin{enumerate}
\item \textbf{Target:} human/prover step validity, cumulative prefix validity,
fixed-completer success probability, or task utility. We recommend intrinsic
step/prefix validity as the primary PRM target and independent answer correctness
as the task metric. Use a separate fixed-completer track for rollout calibration.
\item \textbf{Domain:} a finite candidate pool, trajectories from a frozen
policy, a bounded policy family, or all feature vectors in a known domain.
\item \textbf{Geometry:} independent box, shared parameter ellipsoid,
low-rank bias set, budgeted deviations, or a joint reward distribution.
\item \textbf{Assurance:} empirical coverage, marginal prediction coverage,
candidate-set coverage, uniform function coverage, or a policy improvement bound.
\item \textbf{Use:} robust reranking, trajectory-level RL, robust policy-level
improvement, or task-preserving shaping. These are different objectives.
\end{enumerate}
The methods below are proposals. Their novelty and effectiveness are hypotheses.

\section{Approach A: problem-level calibration of candidate tubes}
\label{sec:approachA}
\textbf{Purpose.} Establish a practical baseline and diagnose whether selection
breaks ordinary coverage. This is an empirical route with a modest, explicit
exchangeability guarantee for a frozen candidate-generation protocol.

\subsection{Construction}
Train a center $\mu(h)\in[0,1]$ and positive scale $s(h)\ge s_{\min}>0$ using
only training problems. Scale candidates include ensemble disagreement,
last-layer leverage, predicted residual magnitude, or quantile widths. Fix the
scorer, generator, parser, and maximum budget $B$ before calibration.
For each calibration problem $i$, generate all $B$ candidates and collect target
labels $y_{ijt}$ for their steps. Problems, not individual steps, are the
exchangeable calibration units. Define
\[
 S_i^-=\max_{j\le B,t\le T_{ij}}\frac{\mu(h_{ijt})-y_{ijt}}{s(h_{ijt})},
 \qquad
 S_i^+=\max_{j\le B,t\le T_{ij}}\frac{y_{ijt}-\mu(h_{ijt})}{s(h_{ijt})}.
\]
On $m$ calibration problems, let $q^\pm$ be the
$\lceil(m+1)(1-\alpha_\pm)\rceil$th order statistic of $S_i^\pm$, augmented
with $+\infty$ if the index exceeds $m$. Replace negative quantiles by zero
to make widths nonnegative. Set $\alpha_-+\alpha_+=\alpha$ and
\[
 L(h)=\max\{0,\mu(h)-q^-s(h)\},\qquad
 U(h)=\min\{1,\mu(h)+q^+s(h)\}.
\]
A one-sided version spends all the error budget on $S^-$ and reports only a
lower prediction bound; it must not be called a two-sided confidence tube.

Under exchangeable problem blocks, independent model fitting, and the same
frozen generation/labeling protocol, split-conformal ranks and a union bound give
\[
 \Prob\{L(h_{jt})\le y_{jt}\le U(h_{jt})\text{ for all candidates and steps
 in a new block}\}\ge1-\alpha.
\]
This adapts standard conformal reasoning \cite{cqr}; it is not a new theorem.
Within-block steps and candidates can be dependent. Any reranker choosing from
that fixed block inherits this containment event. The probability is marginal
over calibration and new problem blocks, not conditional coverage for every
problem or a confidence band over the whole reward function.

\subsection{The target problem, and how to handle it}
If $y$ is a binary validity label, this construction covers the \emph{realized
label}. Many honest intervals may reach zero, making lower-bound selection
uninformative. That is a result to report, not something to fix by relabeling the
interval as epistemic uncertainty. Use sets of valid/invalid labels or abstention
for this track; evaluate informativeness through singleton-set frequency and
false acceptance.

For a success-probability track, keep completer $\pi_c$ fixed and label each
prefix with $K$ independent completion outcomes. Calibration then covers future
rollout estimates $\tilde p$, not $p_{\pi_c}$. If the new block has at most
$M=BT_{\max}$ prefixes, Hoeffding and a union bound give an additional allowance
\[
 \eta_K=\sqrt{\frac{\log(2M/\delta_{\rm MC})}{2K}}.
\]
Expanding each interval by $\eta_K$ covers the true fixed-completer success
probabilities in that block with probability at least
$1-\alpha-\delta_{\rm MC}$. Independent Bernoulli rollouts, a correct answer
checker, and a fixed finite cap are assumptions. This does not convert success
probabilities into intrinsic step correctness.

\subsection{Using the tube}
For the rectangular set and a coordinatewise nondecreasing aggregation $G$,
rerank using
\[
 G(L(h_1),\ldots,L(h_T)).
\]
Repeat for nominal, calibrated-center,
and uncertainty-penalized baselines with the same $G$. A candidate-wise lower
reward may be overly conservative; it should be compared with a shared-function
set in Approach B.

\subsection{Tests, limitations, and novelty}
Sweep nested fixed pools with $N\le B$. Compare stepwise calibration,
trajectory-block calibration, full candidate-block calibration, a union-bound
correction, and a selection-conditional/JOMI adaptation \cite{jomi}. Calibrating
the selected candidate directly requires the selection protocol to be fixed
without circular reuse of its calibration labels.

Problem blocks from an RL-updated generator or a beam search influenced by the
new calibrated reward do not automatically match the frozen-pool guarantee.
Block maxima can be too conservative and annotation cost scales with candidate
count. A publishable result would need a useful improvement in selected-path
coverage and correctness at matched cost, or a clear failure characterization.
The max-score conformal construction itself is standard.

\section{Approach B: a coupled function tube from frozen features}
\label{sec:approachB}
\textbf{Purpose.} Model one shared source of reward error across steps and
policies, providing the cleanest initial theoretical route. AdvPO is close prior
work \cite{advpo}; the proposed extension is process coupling and validation.

\subsection{A tractable model and an exact induced tube}
Start with a frozen feature map $\psi(h)\in\R^d$ and linear real-valued reward
$r_w(h)=\psi(h)^\top w$. Fit a regularized head $\hat w$ and let
\[
 V=\lambda I+\sum_{i=1}^n\psi_i\psi_i^\top,
 \quad \mathcal C=\{w:\norm{w-\hat w}_V\le\beta_n\},
 \quad \U=\{r_w:w\in\mathcal C\}.
\]
Here $\norm{u}_V=\sqrt{u^\top Vu}$. The induced pointwise envelope is
\[
 \psi(h)^\top\hat w\ \pm\ \beta_n\sqrt{\psi(h)^\top V^{-1}\psi(h)}.
\]
Take $\lambda>0$ so that $V$ is invertible. This is a projection of one ellipsoid;
arbitrary simultaneous choices of the
pointwise endpoints are not allowed.

For $y_i=\psi_i^\top w^*+\varepsilon_i$, predictable frozen features,
conditionally $R_\varepsilon$-sub-Gaussian noise, $\norm{w^*}\le S$, and
regularized least squares, a standard confidence choice is
\[
 \beta_n=R_\varepsilon\sqrt{2\log\left(
 \frac{\det(V)^{1/2}}{\det(\lambda I)^{1/2}\delta}\right)}+\sqrt\lambda S.
\]
Self-normalized theory supplies adaptive-data confidence under these assumptions
\cite{abbasi}. This does not justify an unrestricted neural PRM guarantee.
Binary conditional means can fit this model only if the linear mean is correctly
specified. Labels systematically supplied by a biased judge violate it.

For probability outputs $\sigmoid(\psi^\top w)$, monotonicity yields sigmoid
transforms of the logit envelopes \emph{if} the parameter confidence set is
valid. Least-squares confidence theory does not automatically validate a
logistic head. Hessian/Laplace ellipsoids can be empirical baselines; a formal
generalized-linear analysis must establish its own assumptions and radius.

\subsection{Joint geometry and robust scoring}
For a trajectory feature matrix $\Psi_\tau$ whose rows are $\psi(h_t)^\top$,
the induced covariance shape is
\[
 \Sigma_\tau=\Psi_\tau V^{-1}\Psi_\tau^\top.
\]
Repeated steps share uncertainty rather than generating new independent errors.
For a linear aggregation $a^\top\bm r$, where $a$ is fixed for the trajectory,
the exact worst-case score for the unrestricted linear head is
\[
 \inf_{w\in\mathcal C}a^\top\Psi_\tau w
 =a^\top\Psi_\tau\hat w-\beta_n\sqrt{a^\top\Sigma_\tau a}.
\]
This is a support-function identity, not a novel PRM theorem. If rewards must
stay in $[0,1]$, either ensure the ellipsoid satisfies the domain bounds or
intersect it with those constraints and solve the resulting constrained problem.
Clipping outputs or using sigmoid rewards generally destroys the linear identity.
For a minimum objective,
$\inf_w\min_t r_w(h_t)=\min_t\inf_w r_w(h_t)$; for products and nonlinear
means, simultaneous endpoint substitution need not be exact.

Policy-level pessimism differs from per-trajectory pessimism. With
$\bar\psi_\pi=\E_\pi\sum_t a_t\psi(h_t)$, the former equals
\[
 \bar\psi_\pi^\top\hat w-\beta_n\norm{\bar\psi_\pi}_{V^{-1}},
\]
while the latter subtracts
$\beta_n\E_\pi\norm{\sum_t a_t\psi(h_t)}_{V^{-1}}$.
Jensen makes the latter more conservative. Report which objective is optimized.

\subsection{Misspecification and practical construction}
An honest extension is
\[
 r^*(h)=\psi(h)^\top w^*+m(h),\quad |m(h)|\le b_{\rm mis}(h).
\]
The resulting set couples parameter error but allows specified residual error.
Estimate residual scales on separate problems containing adversarial and
on-policy traces. Calling that estimate a uniform bound still requires a
coverage argument. A residual floor can protect against zero ensemble variance,
but should be evaluated through failures rather than assumed effective.

Implement frozen representations, ridge heads, then independently trained
low-rank heads or adapters. Use Cholesky solves rather than explicitly forming
$V^{-1}$. Training needs feature statistics; a full $d\times d$ matrix and
quadratic-form scoring can be expensive. Compare dimension reduction,
diagonal approximations, and low-rank approximations against full geometry.
Do not hide the cost of computing PRM features.

\subsection{Research payoff and risks}
Compare a shared ellipsoid with its rectangular envelope at matched marginal
width and matched joint coverage. Test common-mode bias, repeated steps, and
rare logical transformations. The contribution could be a characterization of
when coupling avoids excessive conservatism or controls policy exploitation.
The main risks are feature misspecification, falsely small radii, and head-only
diversity missing shared backbone errors. A generic last-layer ellipsoid is
already covered by AdvPO and confidence-set literature.

\section{Approach C: structured nuisance and budgeted error tubes}
\textbf{Purpose.} Cover systematic shared biases that a seed ensemble misses.
This route combines empirical attacks with an interpretable robust set.

Let $g(h)$ describe validated nuisance directions: step duplication, verbosity,
unjustified certainty, answer-format markers, or question--solution mismatch.
The last is a corruption detector, not a semantics-preserving nuisance. Define
\[
 r_{u,e}(h)=\hat r(h)-g(h)^\top u-e(h),\quad
 \norm{u}_M\le\rho_b.
\]
For a probability target, intersect these constraints with $0\le r_{u,e}(h)\le1$;
the unconstrained penalties below are then conservative bounds and need not be
the exact constrained optimum. Alternatively, define the set on real-valued
logits and optimize with the explicit link function.
One $u$ applies to the full trajectory or policy; specify that scope. Derive
$g$ from paired transformations whose labels are checked, not arbitrary text
embeddings assumed to represent causality. CoLD motivates the known-bias case
\cite{cold}. Put a zero-centered uncertainty set around unexplained effects;
an estimated known bias should first correct the center.

For nonnegative aggregation weights, an additional budgeted residual model is
\[
 |e_t|\le b_t,\qquad\sum_t\frac{|e_t|}{b_t}\le\Gamma,
 \quad 0\le\Gamma\le T,
\]
omitting zero-width coordinates. Its worst-case linear penalty is the sum of the
largest $\lfloor\Gamma\rfloor$ values of $a_tb_t$, plus the fractional part
times the next largest. This standard robust-set geometry interpolates between
no residual penalty and a rectangular all-steps penalty.

\textbf{Construction protocol.} Create paired correct rewrites and validated
logical corruptions on training problems; fit bias features and residual scales;
estimate a problem-level deviation budget on a disjoint calibration split;
compare an independent box, a shared bias term, and a combined set. If a repeated
systematic mistake affects every step, a small $\Gamma$ is invalid. Use the shared
bias component for such common modes and audit the remaining budget assumptions.

\textbf{Potential contribution.} Show that correlated nuisance uncertainty covers
known failure directions at smaller useful widths than an unstructured box,
then test held-out attacks. Compare against CoLD and hard-negative retraining at
equal label cost. Primary risks are unreliable semantic validation and attacks
outside $g$'s span. Token-count penalties alone are a baseline, not the new method.

\section{Approach D: a joint reward distribution and robust risk}
\textbf{Purpose.} Preserve dependencies in a distributional PRM rather than
attach independent distributions to each step. This is an optional extension
after simpler tube methods work; scalar distributional pessimism is prior work
\cite{unifying,otprm,prism}.

Fit a joint predictive distribution $p(\bm r\mid\tau)$ using shared latent
parameters, low-rank factors, or ensemble trajectories. For a specified $G$,
consider a reward-distribution ambiguity set
\[
 \mathcal Q_\rho(\tau)=\{Q:\KL(Q\|p(\bm r\mid\tau))\le\rho\}.
\]
Its robust expected aggregate is
\[
 \inf_{Q\in\mathcal Q_\rho}\E_Q G(\bm r)
 =\sup_{\eta>0}\{-\eta\log\E_p e^{-G(\bm r)/\eta}-\eta\rho\},
\]
under the standard duality/integrability conditions, with boundary limits
included where needed. The penalized form with fixed $\eta$ omits $-\eta\rho$;
it is a different problem. This variational identity is established robust-risk
machinery, not a novel result.

For real-valued Gaussian rewards and linear $G=a^\top\bm r$, the penalized
score is
$a^\top\bm\mu-(a^\top\Sigma a)/(2\eta)$. Correlation changes the penalty;
dropping off-diagonal terms can severely underestimate shared error. A Gaussian
distribution is not a valid bounded probability distribution unless transformed
or truncated, in which case the simple expression need not hold. Monte Carlo
log-sum-exp estimates must be checked for tail-sampling error.

\textbf{Calibration and experiments.} Train distributional heads on repeated
rollouts in the fixed-completer track or repeated independent annotations in
the validity track, accounting for what those repetitions measure. Compare
independent versus joint models, entropic risk versus quantile bounds, and
conservative distributions versus the same center without penalties. Measure
tail fit after selection and shared errors. Distributional variability alone is
not epistemic confidence; a perfectly fitted wrong label mechanism remains
wrong. Publishability would require a process-specific gain, not re-deriving
scalar entropic rewards.

\section{Approach E: optimize, audit, and update the tube}
\textbf{Purpose.} Make uncertainty estimation respond to the distribution created
by reward optimization. This is the highest-cost combined route.

Define an optimization regime by initial policy, allowed updates, KL/length
caps, algorithm, and budget. A reachable family $\Pi_B$ induces prefix or
trajectory distributions, rather than an unspecified ball around the training
data. Empirically sample a mixture of checkpoints, best-of-$N$ winners, beam
prefixes, and attack-generated traces. A mixture covers its sampled modes;
it does not certify every reachable policy.

\begin{enumerate}
\item Fit the initial PRM and set on training problems. Freeze its parameters
and calibration during a bounded policy-update phase.
\item Optimize the robust score and separately attack it. Select audit examples
with high robust score, low ensemble disagreement despite suspicious reasoning,
and coverage failures; include random on-policy examples to avoid blind sampling.
\item Obtain independent validity labels or fixed-completer rollout targets.
Keep a fresh evaluation stream untouched by reward updates.
\item Retrain the center or uncertainty model, then recalibrate on fresh data.
Compare with equal-budget center-only retraining.
\item Validate new policies after the update; use a predefined accept/defer rule
when evidence of true improvement or valid coverage is insufficient.
\end{enumerate}

If policy $\pi_k$ is fixed before fresh calibration and test blocks are drawn
under it, a per-round guarantee conditional on the previous history can be
combined with $\sum_k\alpha_k\le\alpha$ through a union bound. This only covers
the specified audited deployment blocks. It does not certify policy $\pi_{k+1}$
trained after that calibration, nor reuse of calibration labels to choose among
arbitrary new policies. Uniform function sets or properly justified sequential
inference are needed for broader claims.

Weighted calibration is another comparator when labels have a fixed conditional
mechanism and reference/on-policy density ratios are available \cite{weighted}.
In text, action-product ratios can be extremely variable. Report effective
sample size $(\sum_iw_i)^2/\sum_iw_i^2$, overlap, and clipping bias. Holding the
completer fixed is necessary for the success-probability target; updating the
completer changes that target.

\textbf{Potential contribution.} Measure and repair post-optimization coverage
loss, showing fewer confidently wrong selected trajectories than static sets,
ensembles, and hard-negative training. Main risks are label cost, feedback loops,
catastrophic shared verifier mistakes, and moving targets. Treat an adversarial
curriculum as empirical adaptation, not automatic uniform robustness.

\section{Theoretical characterization: established lemmas and open targets}
\label{sec:theory}
The propositions below are elementary supporting results derived for these notes.
They are not claimed as original publishable theorems. They identify the
assumptions a stronger result must supply.

\begin{proposition}[Rectangular monotone aggregation]
For a finite trajectory and the independent box $\prod_t[L_t,U_t]$, if $G$ is
coordinatewise nondecreasing, then
$\inf_{\bm r\in\prod_t[L_t,U_t]}G(\bm r)=G(\bm L)$.
\end{proposition}
\begin{proof}
Every feasible vector is coordinatewise at least $\bm L$, and $\bm L$ is
feasible. Monotonicity gives both inequalities. For a coupled set $\bm L$ may
not be feasible; its value then remains a conservative lower bound, not an
exact robust score.
\end{proof}

\begin{proposition}[A simultaneous tube protects a selected candidate]
Suppose $R_j^*\in[L_j,U_j]$ for every member of a finite candidate pool.
Let $\hat j\in\arg\max_jL_j$, and let $j^*\in\arg\max_jR_j^*$. Then
\[
 R_{\hat j}^*\ge L_{\hat j},\qquad
 R_{j^*}^*-R_{\hat j}^*\le U_{j^*}-L_{j^*}.
\]
\end{proposition}
\begin{proof}
$R_{\hat j}^*\ge L_{\hat j}\ge L_{j^*}$ and
$R_{j^*}^*\le U_{j^*}$. Subtract. The simultaneous containment event is needed
before selection. For half-widths bounded by $\epsilon$, regret is at most
$2\epsilon$. If these are process-aggregate rewards, the bound concerns that
aggregate and not automatically final-answer correctness.
\end{proof}

\subsection{Why marginal coverage alone fails}
Construct a distribution where a fraction $\alpha$ of candidates are wrong,
receive score one, and receive an interval that excludes the true reward zero.
All other candidates are correct, score $0.8$, and are covered. Random-candidate
coverage is exactly $1-\alpha$. Best-of-$N$ chooses an uncovered wrong candidate
whenever one occurs, with probability $1-(1-\alpha)^N$. At $\alpha=0.05$,
$N=100$, this is about $0.994$. Ties are resolved within wrong candidates.
This is a constructed counterexample, not a measured PRM failure. It illustrates
why selection-aware theory is necessary even before RL.

\begin{proposition}[Robust improvement against a reference]
If the true reward function belongs to a fixed set $\U$, then for every policy
\[
 \Delta_{\U}(\pi):=\inf_{r\in\U}\{J_r(\pi)-J_r(\pi_0)\}
 \le J_{r^*}(\pi)-J_{r^*}(\pi_0).
\]
Therefore $\Delta_{\U}(\pi)>0$ certifies improvement for that target reward.
\end{proposition}
\begin{proof}
Evaluate the infimum at $r^*$. The same containment event works for adaptively
chosen policies if it is a uniform function-containment event on their domain.
\end{proof}
This reference-centered objective is preferable to comparing two separate
pessimistic values: improvement in $\inf_rJ_r(\pi)$ does not alone prove
improvement in the true return. Approximate occupancy estimates, optimization
error, and uncertain labels each require an explicit error allowance.

\subsection{A shift-aware bound with transparent assumptions}
Suppose a bounded trajectory utility satisfies $R^*=\hat R+e$ and
$\E_{P_0}e^2\le\epsilon^2$, with $P_\pi\ll P_0$.
For $w=dP_\pi/dP_0$, Cauchy--Schwarz gives
\[
 \left|\E_{P_\pi}R^*-\E_{P_\pi}\hat R\right|
 \le\epsilon\sqrt{\E_{P_0}w^2}
 =\epsilon\sqrt{1+\chi^2(P_\pi\|P_0)}.
\]
For improvement relative to $P_0$, the error difference is bounded by
$\epsilon\sqrt{\chi^2(P_\pi\|P_0)}$, using $w-1$.
This is a standard change-of-measure calculation related to correlated-proxy
regularization \cite{correlated}; it fails when relevant support is absent or
the ratio is unbounded. Replacing whole trajectories by step occupancy needs
consistent normalization and an additive reward objective.

\subsection{Open theoretical targets}
\begin{enumerate}
\item \textbf{Geometry versus conservatism:} characterize regret or safe
improvement for boxes versus shared ellipsoids at equal joint coverage, including
correlated misspecification and variable-length reasoning. A dual-norm identity
alone is insufficient novelty.
\item \textbf{Coverage under specified optimization:} prove selected-trajectory
coverage for a feasible protocol beyond naive $N T$ union bounds, explicitly
comparing with JOMI and problem-block calibration. Specify the target and budget.
\item \textbf{Necessity and lower bounds:} construct families where any narrow
set calibrated only on a reference policy misses a reachable adversarial mode;
state the needed overlap or complexity assumptions for informative protection.
\item \textbf{Repair-aware process uncertainty:} characterize a set over validity
transitions and show how robust outcome estimation differs from weakest-step
scoring, comparing with RSP. This is more ambitious and needs reliable labels.
\end{enumerate}
A convergence theorem for a neural min--max method should not be asserted from
convex-head or strongly monotone toy assumptions. Prove the restricted model
first, and state the empirical extension separately.

\section{An experiment plan that can support a paper}
\label{sec:experiments}
\subsection{Hypotheses and staged scope}
\begin{description}
\item[H1:] Static pointwise coverage degrades among PRM-selected candidates as
search pressure increases.
\item[H2:] A calibrated tube reduces false high-reward selections better than
equal-cost pessimism, label improvement, or reduced search.
\item[H3:] Coupling improves the coverage--sharpness--accuracy tradeoff when
errors share structure, but can fail under unmodeled common bias.
\item[H4:] Static calibration degrades under policy movement; independently
audited updates recover useful coverage in specified regimes.
\end{description}
Begin with frozen-pool inference and a synthetic setting with known truth.
Proceed to RL only if uncertainty is informative and an overoptimization
baseline can be reproduced. A pure empirical paper can stop at inference-time
claims; it should not imply mitigation of RL hacking without RL experiments.

\subsection{Data and split protocol}
Use PRM800K or appropriate process training data for heads and adapters, with
training, validation, calibration, and test splits separated by problem.
All paraphrases, injected
errors, sampled solutions, and prefixes of a problem belong to one split.
Use ProcessBench, PRMBench, and Socratic-PRMBench as independent diagnostic
benchmarks; do not train a calibrator on their final evaluation labels.

For answer evaluation, generate trajectories on held-out MATH-family problems,
with GSM-style tasks as a second difficulty regime. Small AIME sets are a
stress test with wide uncertainty, not the sole source of statistical claims.
If a published baseline calibrates on MATH500, make that data use explicit and
evaluate on a distinct held-out set. Check upstream training contamination where
possible, and qualify residual uncertainty.

Use at least two generator sizes or families and at least two PRM families,
such as Qwen2.5-Math-PRM and Skywork-o1 PRMs as documented in the cited work.
These are reproducibility candidates, not claims about the latest model releases.
Verify model revisions and inference interfaces before running experiments.

\subsection{Baselines and ablations}
\begin{longtable}{p{0.32\linewidth}p{0.59\linewidth}}
\toprule Comparator & Confound it addresses \\
\midrule\endhead
Nominal single PRM; calibrated center & Is the benefit better point predictions? \\
Mean, minimum, and variance-penalized ensembles & Standard reward pessimism \\
AdvPO-style frozen-head confidence score & Existing confidence-region optimization \\
Quantile PRM; conditional-OT calibration & Existing prefix uncertainty modeling \\
Caution novelty penalty & Single-model OOD conservatism \\
PURE-style minimum; sum/mean/product controls & Credit assignment versus uncertainty \\
Reduced $N$, early stopping, and KL sweep & Less optimization versus a better reward \\
Length correction and hard-negative retraining & Bias correction and improved labels \\
Box versus coupled tube; calibrated versus tuned radius & Geometry and calibration \\
Pointwise, block, union-bound, selection-conditional variants & What survives selection \\
Random widths; constant widths; oracle error widths & Whether width identifies errors \\
Reliable outcome-only RLVR; hybrid rewards & Whether a PRM improves a reliable task signal \\
\bottomrule
\end{longtable}
A full factorial is expensive. Predefine a core set containing nominal,
calibrated-center, ensemble pessimism, AdvPO, one PRM calibration method, Caution,
and the proposed tube; use targeted secondary ablations. Match center quality,
training labels, total generation/scoring cost, and where relevant measured KL.
Radius must be calibrated on calibration data or tuned on validation, never on
the test set. Report those variants separately.

\subsection{Optimization stress tests}
For frozen candidates, sweep $N\in\{1,4,16,64,256\}$, subject to a common token
budget, scoring nested candidate pools. Report both equal-$N$ and equal-compute
comparisons. Test step splitting/merging, repetition, neutral padding, incorrect
numeric transformations, unsupported premises, prompt mismatch, suffix attacks,
and delimiter manipulation. Validate whether each edit preserves or changes
correctness, and keep an untouched transfer family.

For RL, freeze a reward model and tube during each comparison run. Use equal
initial policies, prompts, token caps, training compute, and KL coefficients;
add a matched-observed-KL analysis because penalties can change policy movement.
The primary controlled comparison assigns the same trajectory aggregation to
nominal and robust rewards. Compare process-only and hybrid-outcome supervision
separately. A maximum-empty-response score is a parser defect to fix and report,
not a gain to accept. Save checkpoints and raw rewards before normalization.

Attack the deployed robust score, including its uncertainty head. A defense that
works only against attacks optimized for $\hat r$ may simply move the exploit to
$\hat r-b$. Black-box search and white-box token attacks are separate budgets;
report access assumptions, calls, failed attempts, and transfer results.

\subsection{Metrics and statistical reporting}
The primary metric is held-out answer correctness under increasing pressure,
supplemented by independent process validity on a stratified audit subset.
Use answer checkers with canonicalization and spot checks; do not call a larger
LLM an exact oracle. Score trajectories with correct answers but flawed processes
as a separate category.

Report step, trajectory, candidate-pool, and selected-trajectory coverage;
interval width and singleton/abstention rates; false acceptance of invalid steps;
overestimation error among top-scoring candidates; reward--accuracy curves; length,
repetition, vacuity, KL, and compute. For bounded aligned utilities, also report
lower-bound violations $(L-R^*)_+$ and their tail severity. Do not subtract an
uncalibrated free scalar reward from a binary accuracy metric and interpret the
result as a meaningful numerical ``hacking gap.''

Use paired candidate pools, multiple independent head/policy seeds, and
problem-level bootstrap confidence intervals. A starting plan is at least three
seeds for expensive RL and five for lightweight methods; justify additional runs
from variance and a pilot power analysis. Resampling individual steps as if
independent understates uncertainty. Report all prespecified pressure levels,
including unsuccessful conditions. Specify checkpoints by validation or fixed
training budget, never retrospective test accuracy.

\subsection{Synthetic experiments with known ground truth}
Build finite reasoning trees with known validity labels and optional repair
transitions. Expose features containing a true validity signal and spurious
style dimensions. Vary reference-policy coverage, shared versus independent
reward error, trajectory length, repair rate, and optimization budget. Compare
nominal search/RL, boxes, ellipsoids, bias sets, and oracle bounds. This identifies
when the set really contains truth and when apparent conservatism only shrinks
optimization. Synthetic results validate assumptions, not realistic LLM claims.

\section{Recommended route and publication decision gates}
\begin{longtable}{p{0.15\linewidth}p{0.37\linewidth}p{0.39\linewidth}}
\toprule Route & Minimum convincing result & Main limitation \\
\midrule\endhead
Empirical & Selected-path coverage failures plus a useful, equal-cost repair on
multiple PRMs and pressure regimes & Prediction sets may be too wide; no RL
claim without RL evidence \\
Theoretical & A process-specific coverage/regret or impossibility result beyond
standard confidence-set lemmas & Strong feature/support assumptions may limit
practical implications \\
Combined & A restricted theorem, informative real-PRM tubes, and adaptation-aware
stress tests that match the theorem's mechanism & Highest label and training cost \\
\bottomrule
\end{longtable}

\textbf{Recommended starting point:} Approach A as a diagnostic and calibrated
baseline, Approach B as the structured candidate method, and synthetic tests to
separate shared uncertainty from pointwise errors. Add a misspecification term
from Approach C only when common-mode failures are observed. Approach E is the
RL extension; Approach D is a later alternative, not a prerequisite.

\begin{enumerate}
\item \textbf{Target gate:} a label and reward definition must be coherent.
If success-probability bounds are the only informative ones, narrow the claim to
that target instead of calling them correctness bounds.
\item \textbf{Information gate:} uncertainty must predict overestimation or
useful abstention better than constant/random widths. Otherwise improve the
representation or report why this uncertainty signal fails.
\item \textbf{Mechanism gate:} show gains at fixed aggregation and comparable
compute/optimization pressure. If reduced search explains the effect, do not
claim a stronger reward estimator.
\item \textbf{Prior-work gate:} compare explicitly with AdvPO, PRM quantile/OT
calibration, ensemble pessimism, Caution, and selection-aware inference. Update
the novelty search before drafting a submission.
\item \textbf{Robustness gate:} test adaptive attacks and held-out families. If
calibration fails under RL, retain an honest inference-only result or study the
failure as the main contribution.
\end{enumerate}
The immediate deliverable is a research design, not an experiment result or a
guarantee of publication. A well-supported negative finding about tube coverage
may also be valuable, provided it is broad, controlled, and explains a mechanism
that existing uncertainty methods miss.
